\chapter{導入}
\label{ch:introduction}

Givens–Shortt の論文
\emph{A class of Wasserstein metrics for probability distributions}
（Michigan Math.\ J.\ \textbf{31} (1984), 231–240）の
Proposition 2 を目標とする：
\[
  1\leq p\leq\infty
  \quad\Longrightarrow\quad
  \Wass_p\text{ は }\Pp{p}(X)\text{ 上の距離である．}
\]

\section*{記号の読み替え}

本資料では，原論文の記号を
Villani（\emph{Optimal Transport}, 2009）や
Peyré–Cuturi（\emph{Computational Optimal Transport}, 2019）
に倣い現代的な記法に統一する．

\begin{center}
\renewcommand{\arraystretch}{1.3}
\begin{tabular}{lcc}
  \hline
  & \textbf{原論文} & \textbf{本資料} \\
  \hline
  距離空間       & $(S,\,d)$                      & $(X,\,d)$ \\
  基点           & $0$                            & $x_0$ \\
  確率測度の空間 & $\mathfrak{M}_p(S)$            & $\Pp{p}(X)$ \\
  確率測度       & $P_1,\,P_2$                    & $\mu,\,\nu$ \\
  結合の集合     & $D(P_1,P_2)$                   & $\Couplings(\mu,\nu)$ \\
  Wasserstein 距離 & $W_p(P_1,P_2)$              & $W_p(\mu,\nu)$ \\
  \hline
\end{tabular}
\end{center}
